\documentclass[11pt,a4paper]{article}
\usepackage[utf8]{inputenc}
\usepackage[T1]{fontenc}
\usepackage[swedish]{babel}
\usepackage{amsmath,amssymb}
\usepackage{geometry}
\usepackage{fancyhdr}
\usepackage{array}
\usepackage{booktabs}
\usepackage{xcolor}
\usepackage{hyperref}

\geometry{margin=2.5cm}
\pagestyle{fancy}
\fancyhf{}
\fancyhead[L]{Hvitfeldska Spetsutbildning}
\fancyhead[R]{Matematikhandbok 2026}
\fancyfoot[C]{\thepage}

\definecolor{mainblue}{RGB}{26,84,144}

\title{\Huge\textbf{HVITFELDTSKAS SPETSUTBILDNING}\\[0.5cm]
\Large Komplett Matematikhandbok 2026}
\author{Skapad for Yvonna}
\date{Provdatum: 24 februari 2026}

\begin{document}

\maketitle
\tableofcontents
\newpage

\section{TALTEORI}

\subsection{Delbarhetsregler}

\begin{tabular}{|c|l|l|}
\hline
\textbf{Tal} & \textbf{Regel} & \textbf{Exempel} \\
\hline
2 & Sista siffran ar jamn & 1234 \\
3 & Siffersumman delbar med 3 & 156: 1+5+6=12 \\
4 & Sista tva siffrorna delbara med 4 & 1324: 24/4=6 \\
5 & Slutar pa 0 eller 5 & 1235 \\
6 & Delbar med bade 2 och 3 & 1254 \\
8 & Sista tre siffrorna delbara med 8 & 5016 \\
9 & Siffersumman delbar med 9 & 2286 \\
10 & Slutar pa 0 & 1230 \\
11 & Alternerande siffersumma delbar med 11 & 2728 \\
\hline
\end{tabular}

\subsection{Primtal}

\textbf{Definition:} Ett primtal ar ett heltal $> 1$ som endast ar delbart med 1 och sig sjalv.

\fbox{\parbox{\textwidth}{
\textbf{Primtal under 100:}\\
2, 3, 5, 7, 11, 13, 17, 19, 23, 29, 31, 37, 41, 43, 47, 53, 59, 61, 67, 71, 73, 79, 83, 89, 97
}}

\subsection{Primtalsfaktorisering}

\textbf{Viktiga faktoriseringar:}
\begin{itemize}
\item $36 = 2^2 \times 3^2$
\item $60 = 2^2 \times 3 \times 5$
\item $100 = 2^2 \times 5^2$
\item $180 = 2^2 \times 3^2 \times 5$
\item $350 = 2 \times 5^2 \times 7$
\end{itemize}

\subsection{SGD och MGM}

\textbf{SGD:} Ta gemensamma primfaktorer med \textbf{lagsta} exponent.

\textbf{MGM:} Ta alla primfaktorer med \textbf{hogsta} exponent.

\fbox{$\text{SGD}(a,b) \times \text{MGM}(a,b) = a \times b$}

\subsection{TYPUPPGIFT: Heltalsekvationer}

\textbf{Problem:} $x, y, z$ ar heltal $> 1$. $xy = 350$, $yz = 180$, $xz$ delbar med 8.

\textbf{Losning:}
\begin{enumerate}
\item Faktorisera: $350 = 2 \times 5^2 \times 7$, $180 = 2^2 \times 3^2 \times 5$
\item SGD$(350, 180) = 10$. Mojliga $y$: 2, 5, 10
\item $y = 5$: $x = 70$, $z = 36$, $xz = 2520$ (delbar med 8!)
\end{enumerate}
\textbf{SVAR:} $xz = 2520$

\section{ALGEBRA}

\subsection{Konjugat- och Kvadreringsregler}

\begin{align}
a^2 - b^2 &= (a+b)(a-b)\\
(a+b)^2 &= a^2 + 2ab + b^2\\
(a-b)^2 &= a^2 - 2ab + b^2
\end{align}

\subsection{Potensregler}

\begin{tabular}{|l|l|}
\hline
Multiplikation & $a^m \cdot a^n = a^{m+n}$ \\
Division & $a^m / a^n = a^{m-n}$ \\
Potens av potens & $(a^m)^n = a^{mn}$ \\
Negativ exponent & $a^{-n} = 1/a^n$ \\
Nollexponent & $a^0 = 1$ \\
\hline
\end{tabular}

\subsection{Andragradsekvationer}

\textbf{PQ-formeln:} For $x^2 + px + q = 0$:
\[ x = -\frac{p}{2} \pm \sqrt{\left(\frac{p}{2}\right)^2 - q} \]

\textbf{ABC-formeln:} For $ax^2 + bx + c = 0$:
\[ x = \frac{-b \pm \sqrt{b^2 - 4ac}}{2a} \]

\subsection{Talföljder}

\textbf{Aritmetisk:}
$a_n = a_1 + (n-1)d$, \quad $S_n = \frac{n(a_1 + a_n)}{2}$

\textbf{Geometrisk:}
$a_n = a_1 \cdot k^{n-1}$, \quad $S_n = a_1 \cdot \frac{k^n - 1}{k - 1}$

\section{GEOMETRI}

\subsection{Trianglar}

\textbf{Vinkelsumma:} $180°$

\fbox{\textbf{Pythagoras:} $a^2 + b^2 = c^2$}

\textbf{Pythagoriska tripplar:} 3-4-5, 5-12-13, 8-15-17, 7-24-25

\textbf{Area:} $A = \frac{bh}{2}$

\subsection{Speciella trianglar}

\textbf{Liksidig} (sida $a$): Hojd $h = \frac{a\sqrt{3}}{2}$, Area $A = \frac{a^2\sqrt{3}}{4}$

\textbf{30-60-90:} Sidor $1 : \sqrt{3} : 2$

\textbf{45-45-90:} Sidor $1 : 1 : \sqrt{2}$

\subsection{Cirkeln}

Omkrets: $O = 2\pi r$, \quad Area: $A = \pi r^2$

Baglangd: $b = \frac{v}{360} \cdot 2\pi r$, \quad Sektorarea: $A = \frac{v}{360} \cdot \pi r^2$

\textbf{Viktigt:} Tangent ar vinkelrat mot radie!

\subsection{Koordinatgeometri}

Avstand: $d = \sqrt{(x_2-x_1)^2 + (y_2-y_1)^2}$

Mittpunkt: $M = \left(\frac{x_1+x_2}{2}, \frac{y_1+y_2}{2}\right)$

Lutning: $k = \frac{y_2 - y_1}{x_2 - x_1}$

Vinkelrata linjer: $k_1 \cdot k_2 = -1$

\section{KOMBINATORIK}

\subsection{Grundprinciper}

\textbf{Multiplikation:} $n_1 \times n_2 \times \ldots \times n_k$ satt

\textbf{Komplement:} Gynnsamma = Totalt $-$ Ogynnsamma

\subsection{Formler}

Fakultet: $n! = n \times (n-1) \times \ldots \times 1$

Permutation: $P(n,r) = \frac{n!}{(n-r)!}$

Kombination: $\binom{n}{r} = \frac{n!}{r!(n-r)!}$

\subsection{Vagrakning}

$m$ steg hoger + $n$ steg uppat = $\binom{m+n}{m}$ vagar

Exempel: $5 \times 5$ rutnat: $\binom{10}{5} = 252$ vagar

\section{SANNOLIKHET}

\[ P(A) = \frac{\text{gynnsamma}}{\text{mojliga}} \]

Komplement: $P(A') = 1 - P(A)$

\textbf{Minnesregel:} P(minst en) = 1 $-$ P(ingen)

\section{ORDPROBLEM}

\textbf{Samarbete:} $t = \frac{ab}{a+b}$

\textbf{Motas:} $t = \frac{\text{stracka}}{v_1 + v_2}$

\textbf{Ikapp:} $t = \frac{\text{forsprang}}{v_{\text{snabb}} - v_{\text{langsam}}}$

\section{MINNESREGLER}

\begin{itemize}
\item \textbf{3 och 9:} Siffersumma!
\item \textbf{SGD/MGM:} Minsta/Storsta exponenter
\item \textbf{Kombinatorik:} Ordning = Permutation, Annars = Kombination
\item \textbf{Vagrakning:} = Kombination!
\item \textbf{Tangent:} Vinkelrat mot radie
\end{itemize}

\vfill
\begin{center}
\rule{10cm}{0.5pt}\\[0.5cm]
{\Large\textbf{Lycka till, Yvonna!}}\\
Version 1.0 -- December 2024
\end{center}

\end{document}
